% module_interfaces_tables.tex
% Interface (port) tables for the 11 compute-core modules of the ECG CNN accelerator.
% Widths and signal names taken verbatim from hardware/RTL/*.v (NOT RTL/txt/,
% NOT hardware/RTL_weight/ — the production weight-load + runtime-topology-config
% variant has extra cfg_*/w_wr_*/b_wr_*/fcw_wr_*/out_ch_mask ports not present here).
% Paste-ready booktabs fragments — \input this file. Needs: \usepackage{booktabs,array}
% Companion Markdown master: module_interfaces.md
%
% Column convention: Dir | Width | Name | Group | Purpose.
% "s" prefix on Dir (in.s / out.s) marks a signed port.

% =====================================================================
\begin{table}[t]
\centering
\caption{Interface of \texttt{cp\_mac} — MAC datapath (S1--S4).}
\label{tab:if-cp-mac}
\small
\setlength{\tabcolsep}{4pt}
\renewcommand{\arraystretch}{1.1}
\begin{tabular}{@{}llll@{}}
\toprule
\textbf{Dir} & \textbf{Width} & \textbf{Name} & \textbf{Purpose} \\
\midrule
in    & 1        & \texttt{clk}      & clock \\
in    & [39:0]   & \texttt{x\_in}    & 5 samples, packed 5$\times$8b \\
in    & [39:0]   & \texttt{w}        & 5 weights, packed 5$\times$8b \\
out.s & [19:0]   & \texttt{tree\_out}& adder-tree output ($\Sigma$ 5 products) \\
\bottomrule
\end{tabular}
\end{table}

% =====================================================================
\begin{table}[t]
\centering
\caption{Interface of \texttt{cp\_accumulate\_rescale} — accumulate + bias + rescale + ReLU (S5--S8).}
\label{tab:if-cp-accres}
\small
\setlength{\tabcolsep}{4pt}
\renewcommand{\arraystretch}{1.1}
\begin{tabular}{@{}llll@{}}
\toprule
\textbf{Dir} & \textbf{Width} & \textbf{Name} & \textbf{Purpose} \\
\midrule
in    & 1      & \texttt{clk}            & clock \\
in    & 1      & \texttt{rst}            & reset \\
in    & 1      & \texttt{pool\_rst}      & pool/layer-transition reset \\
in.s  & [19:0] & \texttt{tree\_out}      & MAC output (from cp\_mac) \\
in.s  & [31:0] & \texttt{bias\_in}       & INT32 bias (from cp\_weight\_store) \\
in    & [3:0]  & \texttt{a\_in}          & channel counter, delay 5 cy (a\_d5) \\
in    & [3:0]  & \texttt{in\_ch}         & IN\_CH of layer (1/4/4/8) \\
in    & 1      & \texttt{compute\_en\_in}& pipeline enable, delay 5 cy (ce\_d5) \\
in    & [3:0]  & \texttt{nb}             & rescale shift (0..15; max 8) \\
in    & 1      & \texttt{relu\_en}       & 1 = Conv4 only \\
out.s & [7:0]  & \texttt{relu\_out}      & INT8 activation \\
out   & 1      & \texttt{relu\_v}        & valid \\
\bottomrule
\end{tabular}
\end{table}

% NOTE: internal stage S\_bias (register \texttt{biased}/\texttt{bias\_valid},
% gated by \texttt{out\_valid\_d1}) is a bare delay-then-capture with no adder —
% it replaces an older two-stage S5b (\texttt{acc\_final\_r}/\texttt{acc\_final\_v})
% + S\_bias pair (commit 369f200, "fold S5b into S\_bias"). Port list above is
% unaffected (internal-only change); see cp\_pipeline.md / cp\_submodule\_timing.md.

% =====================================================================
\begin{table}[t]
\centering
\caption{Interface of \texttt{cp\_pool} — MaxPool rolling comparator (S9).}
\label{tab:if-cp-pool}
\small
\setlength{\tabcolsep}{4pt}
\renewcommand{\arraystretch}{1.1}
\begin{tabular}{@{}llll@{}}
\toprule
\textbf{Dir} & \textbf{Width} & \textbf{Name} & \textbf{Purpose} \\
\midrule
in    & 1     & \texttt{clk}            & clock \\
in    & 1     & \texttt{rst}            & reset \\
in    & 1     & \texttt{pool\_rst}      & pool/layer reset \\
in.s  & [7:0] & \texttt{relu\_out}      & INT8 activation \\
in    & 1     & \texttt{relu\_v}        & valid \\
in    & 1     & \texttt{compute\_en\_in}& original gate (ce\_d5) \\
out   & 1     & \texttt{pool\_write}    & write strobe (AND cp\_en externally) \\
out.s & [7:0] & \texttt{pool\_out}      & max-pooled value to Pong SRAM \\
\bottomrule
\end{tabular}
\end{table}

% =====================================================================
\begin{table}[t]
\centering
\caption{Interface of \texttt{cp\_block} — Conv-Pool block, 1 output channel
(\texttt{IN\_CH\_W}$=4$).}
\label{tab:if-cp-block}
\small
\setlength{\tabcolsep}{4pt}
\renewcommand{\arraystretch}{1.1}
\begin{tabular}{@{}llll@{}}
\toprule
\textbf{Dir} & \textbf{Width} & \textbf{Name} & \textbf{Purpose} \\
\midrule
in    & 1      & \texttt{clk}            & clock \\
in    & 1      & \texttt{rst}            & reset \\
in    & [39:0] & \texttt{x\_in}          & 5 taps (mux\_s1), packed 5$\times$8b \\
in    & [39:0] & \texttt{w}              & weights (w\_packed), packed 5$\times$8b \\
in.s  & [31:0] & \texttt{bias\_in}       & INT32 bias (from cp\_weight\_store) \\
in    & [3:0]  & \texttt{a\_in}          & channel counter, delay 5 cy (a\_d5) \\
in    & [3:0]  & \texttt{in\_ch}         & IN\_CH of layer (1/4/4/8) \\
in    & 1      & \texttt{compute\_en\_in}& pipeline enable, delay 5 cy (ce\_d5) \\
in    & [3:0]  & \texttt{nb}             & rescale shift (0..15; max 8) \\
in    & 1      & \texttt{relu\_en}       & 1 = Conv4 only \\
in    & 1      & \texttt{pool\_rst}      & pool reset (layer transition) \\
out   & 1      & \texttt{pool\_write}    & write strobe (AND cp\_en externally) \\
out.s & [7:0]  & \texttt{pool\_out}      & max-pooled value to Pong SRAM \\
\bottomrule
\end{tabular}
\end{table}

% Pipeline: MULT(1) -> TREE(3) -> ACC(IN_CH) -> S_bias(1, bare capture, no adder)
% -> RESCALE(2) -> RELU(1) -> POOL. (S_bias replaces the old ACC_FINAL(1)+BIAS(1) pair.)

% =====================================================================
\begin{table}[t]
\centering
\caption{Interface of \texttt{cp\_weight\_store} — weight + bias storage
(ROM single-load build: 4 per-layer FF-array ROMs, no bus write).}
\label{tab:if-cp-wstore}
\small
\setlength{\tabcolsep}{4pt}
\renewcommand{\arraystretch}{1.1}
\begin{tabular}{@{}llll@{}}
\toprule
\textbf{Dir} & \textbf{Width} & \textbf{Name} & \textbf{Purpose} \\
\midrule
in  & 1       & \texttt{clk}           & clock \\
in  & 1       & \texttt{rst}           & reset (unused in logic; kept for port symmetry) \\
in  & [2:0]   & \texttt{layer\_state}  & CONV1=2 .. CONV4=5 \\
in  & [3:0]   & \texttt{a}             & input-channel counter (word index) \\
out & [319:0] & \texttt{w\_packed\_flat}& 8$\times$40b weights (reg N+1) \\
out & [255:0] & \texttt{b\_cur\_flat}  & 8$\times$32b INT32 bias (current layer) \\
\bottomrule
\end{tabular}
\end{table}

% NOTE: no cfg_base / w_wr_* / b_wr_* ports in this (RTL/) build — weights are
% baked in via $readmemh (conv1_w..conv4_w.hex, conv_bias.hex) at elaboration.
% The bus-writable M10K weight-RAM variant with cfg_base lives in RTL_weight/.

% =====================================================================
\begin{table}[t]
\centering
\caption{Interface of \texttt{cp\_engine} — 8 parallel CP blocks + SRW + address gen.}
\label{tab:if-cp-engine}
\small
\setlength{\tabcolsep}{4pt}
\renewcommand{\arraystretch}{1.1}
\begin{tabular}{@{}llll@{}}
\toprule
\textbf{Dir} & \textbf{Width} & \textbf{Name} & \textbf{Purpose} \\
\midrule
in  & 1      & \texttt{clk}              & clock \\
in  & 1      & \texttt{rst}              & reset \\
in  & [3:0]  & \texttt{a}                & channel counter 0..IN\_CH$-$1 \\
in  & [3:0]  & \texttt{in\_ch}           & IN\_CH: 1/4/4/8 \\
in  & [11:0] & \texttt{in\_len}          & IN\_LEN (2500/500/100/20) \\
in  & 1      & \texttt{shift\_en}        & = (a == IN\_CH$-$1) \\
in  & 1      & \texttt{srw\_rst}         & SRW clear pulse \\
in  & 1      & \texttt{compute\_en}      & pipeline enable (0 in pre-fetch) \\
in  & [3:0]  & \texttt{nb}               & rescale shift/layer (max 8) \\
in  & 1      & \texttt{relu\_en}         & 1 = Conv4 only \\
in  & [7:0]  & \texttt{cp\_en}           & active output-channel bitmask \\
in  & [2:0]  & \texttt{layer\_state}     & CONV1=2 .. CONV4=5 \\
in  & 1      & \texttt{pool\_rst}        & reset pool\_cnt (layer transition) \\
in  & [7:0]  & \texttt{input\_sram\_dout}& from input\_sram (Conv1 only) \\
in  & [63:0] & \texttt{ping\_dout}       & from ping\_pong: [ch*8+:8] \\
out & [63:0] & \texttt{pong\_din}        & write data/channel: [ch*8+:8] \\
out & [7:0]  & \texttt{pong\_we}         & per-channel write enable \\
out & [11:0] & \texttt{sram\_rd\_addr}   & read addr $\to$ input/ping\_pong \\
in  & [11:0] & \texttt{sram\_rd\_addr\_in}& base from controller (= t) \\
\bottomrule
\end{tabular}
\end{table}

% NOTE: no w_wr_*/b_wr_*/cfg_base ports in this (RTL/) build — those belong to
% the RTL_weight/ production variant (weight-RAM bus write + runtime layer base).

% =====================================================================
\begin{table}[t]
\centering
\caption{Interface of \texttt{cnn\_controller} — unified FSM (topology hard-coded).}
\label{tab:if-cnn-ctrl}
\small
\setlength{\tabcolsep}{4pt}
\renewcommand{\arraystretch}{1.1}
\begin{tabular}{@{}llll@{}}
\toprule
\textbf{Dir} & \textbf{Width} & \textbf{Name} & \textbf{Purpose} \\
\midrule
in  & 1      & \texttt{clk}            & clock \\
in  & 1      & \texttt{rst}            & synchronous reset \\
in  & 1      & \texttt{start}          & 1-cycle start pulse \\
in  & 1      & \texttt{pool\_write}    & from cp\_engine (ch0 representative) \\
out & [3:0]  & \texttt{a}              & channel counter 0..IN\_CH$-$1 \\
out & [11:0] & \texttt{t}              & output position counter \\
out & 1      & \texttt{shift\_en}      & = (a == in\_ch$-$1) \\
out & 1      & \texttt{srw\_rst}       & SRW clear pulse (1 cy) \\
out & 1      & \texttt{compute\_en}    & pipeline enable \\
out & [3:0]  & \texttt{in\_ch}         & current layer IN\_CH \\
out & [11:0] & \texttt{in\_len}        & current layer IN\_LEN \\
out & [3:0]  & \texttt{nb}             & rescale shift (max 8) \\
out & 1      & \texttt{relu\_en}       & ReLU enable (Conv4 only) \\
out & [7:0]  & \texttt{cp\_en}         & active output-channel bitmask \\
out & 1      & \texttt{bank\_sel}      & Ping/Pong bank selector \\
out & [11:0] & \texttt{pong\_addr}     & Pong SRAM write address \\
out & 1      & \texttt{pool\_rst}      & reset pool\_cnt (layer transition) \\
out & [2:0]  & \texttt{fc\_sub\_state} & GAP/FC/flush/argmax/done sub-state \\
out & [3:0]  & \texttt{gap\_step}      & 0..5 \\
out & [3:0]  & \texttt{fc\_step}       & 0..9 \\
out & [1:0]  & \texttt{argmax\_step}   & 0..3 \\
in  & [1:0]  & \texttt{argmax\_result} & argmax class (from gap\_fc) \\
out & [2:0]  & \texttt{layer\_state}   & exposed for cp\_engine MUX + top \\
out & 1      & \texttt{busy}           & 1 while $\neq$ IDLE/DONE \\
out & 1      & \texttt{done}           & 1-cycle pulse \\
out & [1:0]  & \texttt{result}         & latched argmax class \\
\bottomrule
\end{tabular}
\end{table}

% NOTE: no cfg_in_ch/cfg_cp_en/cfg_nb ports in this (RTL/) build. Topology is
% hard-coded via 3 internal Verilog functions (cfg_in_ch_of/cfg_cp_en_of/cfg_nb_of,
% indexed by layer 0..3, NOT bus ports): in_ch=1,4,4,8; cp_en=0F,0F,FF,FF;
% nb=8,7,6,7 (ningbo retrain 2026-07-28; Chapman-only build used nb2=6).
% Runtime-loadable cfg_* ports belong to RTL_weight/.

% =====================================================================
\begin{table}[t]
\centering
\caption{Interface of \texttt{gap\_fc\_argmax} — GAP $\to$ FC $\to$ Argmax
(thin wrapper over gap\_unit / fc\_unit / argmax\_unit).}
\label{tab:if-gap-fc}
\small
\setlength{\tabcolsep}{4pt}
\renewcommand{\arraystretch}{1.1}
\begin{tabular}{@{}llll@{}}
\toprule
\textbf{Dir} & \textbf{Width} & \textbf{Name} & \textbf{Purpose} \\
\midrule
in  & 1      & \texttt{clk}           & clock \\
in  & 1      & \texttt{rst}           & reset \\
in  & [2:0]  & \texttt{fc\_sub\_state}& GAP/FC/flush/argmax/done \\
in  & [3:0]  & \texttt{gap\_step}     & 0..5 \\
in  & [3:0]  & \texttt{fc\_step}      & 0..9 \\
in  & [1:0]  & \texttt{argmax\_step}  & 0..3 \\
in  & [63:0] & \texttt{ping\_dout}    & Conv4 output: [ch*8+:8], 1-cy \\
out & [8:0]  & \texttt{gap\_rd\_addr} & Ping read addr, broadcast 8 ch (0..3) \\
out & [1:0]  & \texttt{result}        & argmax class index \\
\bottomrule
\end{tabular}
\end{table}

% NOTE: no out_ch_mask/fcw_wr_* ports in this (RTL/) build — GAP always treats
% all 8 Conv4 channels as active (fixed topology); FC weight/bias are baked in
% via $readmemh (fc_weights.hex/fc_bias.hex) inside fc_unit, no bus write.
% The Conv4-out_ch<8 (out_ch_mask) and FC bus-load path belong to RTL_weight/.

% =====================================================================
\begin{table}[t]
\centering
\caption{Interface of \texttt{input\_sram} — 2500$\times$8b simple dual-port (M10K).}
\label{tab:if-input-sram}
\small
\setlength{\tabcolsep}{4pt}
\renewcommand{\arraystretch}{1.1}
\begin{tabular}{@{}llll@{}}
\toprule
\textbf{Dir} & \textbf{Width} & \textbf{Name} & \textbf{Purpose} \\
\midrule
in  & 1      & \texttt{clk}      & clock \\
in  & [11:0] & \texttt{wr\_addr} & write addr 0..2499 \\
in  & [7:0]  & \texttt{din}      & write data \\
in  & 1      & \texttt{we}       & write enable \\
in  & [11:0] & \texttt{rd\_addr} & read addr 0..2499 \\
out & [7:0]  & \texttt{dout}     & read data (1-cycle sync) \\
\bottomrule
\end{tabular}
\end{table}

% =====================================================================
\begin{table}[t]
\centering
\caption{Interface of \texttt{ping\_pong\_sram} — inter-layer feature-map buffer
(2 sets $\times$ 8 ch $\times$ 500 $\times$ 8b).}
\label{tab:if-pingpong}
\small
\setlength{\tabcolsep}{4pt}
\renewcommand{\arraystretch}{1.1}
\begin{tabular}{@{}llll@{}}
\toprule
\textbf{Dir} & \textbf{Width} & \textbf{Name} & \textbf{Purpose} \\
\midrule
in  & 1      & \texttt{clk}      & clock \\
in  & 1      & \texttt{bank\_sel}& 0: A=Ping B=Pong $\mid$ 1: swapped \\
in  & [8:0]  & \texttt{wr\_addr} & Pong write addr 0..499 \\
in  & [63:0] & \texttt{din}      & 8 ch packed: [ch*8+:8] \\
in  & [7:0]  & \texttt{we}       & per-channel write enable \\
in  & [8:0]  & \texttt{rd\_addr} & Ping read addr 0..499 \\
out & [63:0] & \texttt{dout}     & 8 ch packed: [ch*8+:8] \\
\bottomrule
\end{tabular}
\end{table}

% =====================================================================
\begin{table}[t]
\centering
\caption{Interface of \texttt{ecg\_core} — bus-agnostic compute core
(pp + cpe + gfa + ctrl; \texttt{input\_sram} lives in the wrapper).}
\label{tab:if-ecg-core}
\small
\setlength{\tabcolsep}{4pt}
\renewcommand{\arraystretch}{1.1}
\begin{tabular}{@{}llll@{}}
\toprule
\textbf{Dir} & \textbf{Width} & \textbf{Name} & \textbf{Purpose} \\
\midrule
in  & 1      & \texttt{clk}            & clock \\
in  & 1      & \texttt{rst}            & synchronous reset (active high) \\
out & [11:0] & \texttt{input\_rd\_addr}& read addr $\to$ wrapper input\_sram \\
in  & [7:0]  & \texttt{input\_dout}    & read data (1-cy latency) \\
in  & 1      & \texttt{start}          & kick off inference \\
out & 1      & \texttt{busy}           & busy \\
out & 1      & \texttt{done}           & done pulse \\
out & [1:0]  & \texttt{result}         & class 0..3 \\
\bottomrule
\end{tabular}
\end{table}

% NOTE: this (RTL/) ecg_core has only the 7 ports above — a minimal bus-agnostic
% core. The RTL_weight/ production variant adds w_wr_*/b_wr_*/fcw_wr_* (weight/
% bias/FC load) and cfg_in_ch/cfg_cp_en/cfg_nb/cfg_base (runtime topology); its
% gap_fc_argmax out_ch_mask is wired from cfg_cp_en[3*8 +: 8] (Conv4 slice)
% inside ecg_core.

% =====================================================================
% Build variants: RTL/ has NONE (single ROM build, no ifdef/define). The
% WEIGHT_ROM / NO_WEIGHT_INIT defines that select between FF-ROM and M10K
% weight-RAM reload apply only to RTL_weight/cp_weight_store.v.
